\begin{tikzpicture}[font=\scriptsize,cell/.style={minimum width=4.2mm,minimum height=4.2mm,inner sep=0pt,draw=hmtgray!35}]
  \begin{scope}[xshift=0cm]
    \node[font=\bfseries\scriptsize,hmtblue,align=center] at (1.72,.92)
      {matriz aritmética APP\\[-.5mm]\(\Pi(i,j)=\dr(ij)\)};
    \foreach \i in {1,...,9}{
      \node[font=\bfseries] at (-.35,{-(\i-1)*.43}) {\i};
      \foreach \j in {1,...,9}{
        \pgfmathtruncatemacro{\v}{mod(\i*\j-1,9)+1}
        \pgfmathtruncatemacro{\border}{(\i==9)||(\j==9)}
        \pgfmathtruncatemacro{\rad}{(mod(\i,3)==0)||(mod(\j,3)==0)}
        \ifnum\border=1
          \node[cell,fill=hmtlightblue,text=hmtblue,font=\bfseries] at ({(\j-1)*.43},{-(\i-1)*.43}) {\v};
        \else
          \ifnum\rad=1
            \node[cell,fill=hmtlightviolet] at ({(\j-1)*.43},{-(\i-1)*.43}) {\v};
          \else
            \node[cell,fill=white] at ({(\j-1)*.43},{-(\i-1)*.43}) {\v};
          \fi
        \fi
      }
    }
    \foreach \j in {1,...,9}{\node[font=\bfseries] at ({(\j-1)*.43},.38) {\j};}
    \draw[hmtgold,very thick,rounded corners=1pt] (1.075,-1.075) rectangle (1.935,-1.935);
    \draw[hmtblue,very thick] (3.72,.22)--(3.72,-3.72)--(-.22,-3.72);
    \node[hmtblue,anchor=north] at (1.75,-4.02) {subconjunto coordenado de la clase nula};
  \end{scope}

  \begin{scope}[xshift=5.25cm]
    \node[draw=hmtblue,rounded corners=2pt,fill=hmtlightblue,
      text width=6.25cm,align=left,inner sep=8pt] at (2.9,-1.35) {%
      \textbf{Ley interna de acumulación}\par\smallskip
      \[
        \Pi(i,j)=\dr(\underbrace{i+\cdots+i}_{j\ {\rm sumandos}})
        =\dr(ij).
      \]
      La suma genera las filas multiplicativas y, en la dirección
      transversal, las columnas. El observable
      \(\Sigma(i,j)=\dr(i+j)\) conserva esa ley elemental para comparar
      ambas lecturas sobre el mismo punto.};
    \node[draw=hmtgold,rounded corners=2pt,fill=hmtlightgold,
      text width=6.25cm,align=center,inner sep=8pt] at (2.9,-3.65) {%
      \(\displaystyle
      \AAPP=(X_9,\GraphAPP,\Pi;\Sigma,
      \operatorname{Path}(\GraphAPP))\)\par
      una sola APP, una matriz principal, una ley aditiva interna y cuatro
      direcciones orientadas \(N,E,S,O\).};
  \end{scope}
\end{tikzpicture}
